\begin{papercard}{hayes2007}
\cf{Stages} \stagetag{I}. A short methods note proposing Krippendorff's $\alpha$ as the \emph{standard} reliability measure for coded data, with five criteria for a good measure, a one-paragraph verdict on each competing index, and a worked example with the KALPHA macro (SPSS/SAS) including bootstrapped confidence intervals. \emph{Note:} the appended file is the author manuscript. Its text layer prints $\alpha$ as ``a'', and the quotations below follow that layer.

\cf{Problem} Coded data can be trusted only after reliability is shown, but many indices exist (Popping compares 43 for nominal data, mostly for two observers) and respond to different properties \pc{hayes2007}{2}{Choosing an index of reliability is complicated by the number of indices that have been proposed}. Lombard et al.\ called for a standard but did not name one. Of the three kinds of reliability (stability, reproducibility, accuracy), reproducibility is ``arguably the strongest and most feasible kind to test'' \pc{hayes2007}{2}{reproducibility arguably is the strongest and most feasible kind to test}.

\cf{Preconditions for reliability data} Units must be independent, coders must work independently from common instructions, the reliability units should be a random sample of the data in question, and observers should be common enough to be found elsewhere \pc{hayes2007}{3}{Before any statistic can measure reliability, it must be applied to reliability data properly obtained}.

\cf{Five criteria for a good index} (1) Assess agreement among two or more observers, independent of their number and invariant to permutation and selective participation. (2) Be grounded in the distribution of categories actually used, not confounded by the number of categories available. (3) Have a numerical scale between two points with sensible interpretations: perfect agreement (1) and the absence of any statistical relationship between units and their descriptions (0). (4) Fit the level of measurement and keep the same mathematical structure across metrics. (5) Have a known or computable sampling behaviour \pc{hayes2007}{3}{Given such reliability data, a good index of reliability should have the following properties}.

\cf{Verdicts on competing indices} \textbf{Percent agreement} fails criteria 2 and 3 (0\,\% is unattainable with independent coders) \pc{hayes2007}{3}{This measure is easily calculated but flawed in nearly all important respects}. \textbf{Bennett's $S$} (reinvented as Guilford's G, RE, $C$, $k_n$, $I_r$) is inflated by unused categories \pc{hayes2007}{4}{S is inflated by the number of unused categories that the author of the instrument had imagined}. \textbf{Scott's $\pi$} was the first to satisfy 2 and 3 but is limited to two coders and nominal data. \textbf{Cohen's $\kappa$} ``created a hybrid index'': it violates (1) by making observers non-interchangeable and (3) by a zero point suited to association statistics, ``punishing'' coders for agreeing on category frequencies and letting systematic disagreements inflate $\kappa$ \pc{hayes2007}{4}{Kappa violates (1) by not allowing observers to be freely permutable or interchangeable}. \textbf{Fleiss' $K$} generalises $\pi$ to many observers, but only for nominal data. \textbf{Cronbach's $\alpha$} measures consistency, not agreement \pc{hayes2007}{4}{Cronbach’s (1951) alpha (aC) is a statistic for interval-level data that responds to the consistency of observers}.

\cf{Krippendorff's $\alpha$} Satisfies all five: it counts pairs of values assigned to units, unaffected by the number of observers (``dispels the common belief that reliability is the more difficult the more observers are involved''), roots expectations in the data, anchors 1 and 0, handles nominal to ratio data and \emph{missing data}, and its sampling distribution is bootstrapped \pc{hayes2007}{5}{Krippendorff’s (1970, 2004a) alpha (a) satisfies all of the above conditions}. Special cases: nominal two-observer $\alpha$ is asymptotically Scott's $\pi$; ordinal equals Spearman's $\rho$; interval equals the intraclass correlation \pc{hayes2007}{5}{In its two-observer nominal data version, a is asymptotically equal to Scott’s p}.

\cf{Worked example (KALPHA)} Five observers rated the tone of 40 randomly drawn news articles (0--3 scale, with missing judgments), after training and independent reading. Command: \texttt{KALPHA judges = obs1 \ldots\ obs5/level = 2/detail = 1/boot = 10000} (level 1--4 = nominal to ratio) \pc{hayes2007}{6}{The data for this example come from five observers who were asked to evaluate the local news coverage}. Result: ordinal $\alpha=.7598$, 95\,\% CI [.7078, .8078], 239 pairs. Probability of failing $\alpha_{\min}$: $q=.9473$ for .80, $q=.0125$ for .70. Treated as nominal $\alpha=.4765$, interval .7574, ratio .6621, showing that the disagreements follow the ordinal metric \pc{hayes2007}{7}{In other words, if this research problem demands that reliability must not be below amin = 0.800}. Footnote: the null hypothesis ``agreement above chance'' is irrelevant, since data must not deviate significantly from \emph{perfect} reliability \pc{hayes2007}{6}{However, this null-hypothesis is irrelevant where data reliability is considered}. Bootstrapping resamples the \emph{pairs of judgments} per unit; about 10,000 resamples suffice \pc{hayes2007}{7}{The bootstrap sampling distribution of alpha is generated by taking a random sample of 239 pairs of judgments}.

\cf{Use for this thesis} Practical justification for choosing $\alpha$ over $\kappa$ and a how-to for reporting it: state the metric (nominal for categorical codes, ordinal for graded answers), the number of units and coders, the missing-data pattern, $\alpha$ with a bootstrapped CI, and $q$ against a pre-stated $\alpha_{\min}$. KALPHA is SPSS/SAS, so use an R or Python implementation (check it against a hand calculation) and record the version. It is an advocacy paper by the coefficient's author, so cite the criticisms of $\kappa$ and AC1 as a position, and balance it against Hallgren, McHugh and Gwet. For heavily skewed binary codes see the low-variation caveat in Krippendorff (2019).
\end{papercard}
